% GENERATED FILE. Edit canonical lessons, then run npm run generate:books.
% canonical-source-hash: fnv1a64:dde846dba2ee7a13
% canonical-lessons: MR-C481-padse, MR-C481-kaph, MR-C481-sinkne, MR-C481-sinkarne, MR-C481-khakarne, MR-R481-review-474-481

\chapter{Coughs and Colds}
\label{ch:mr-coughs-and-colds}
\hlchaptermodality{481}

\begin{chapteropening}
\textbf{By the end of this chapter:} \emph{I can talk about a runny cold, sneezing and clearing my throat.}

Complete the last lesson of chapter 481: I can talk about a runny cold, sneezing and clearing my throat.
\end{chapteropening}

\section[\texorpdfstring{\mr{पडसे} (paḍse)}{पडसे (paḍse)}]{\mr{पडसे} (paḍse) — a runny cold}

\label{lesson:MR-C481-padse}



\begin{quote}
Before the new one: say the Marathi for weak, then the Marathi for to feel sick.
\end{quote}

\subsection*{You'll want to know: \mr{पडसे}}
\textbf{\mr{पडसे}} — \emph{paḍse} — \textquotedblleft{}a runny cold\textquotedblright{}.

\begin{quote}
\mr{हवा} \mr{बदलली} \mr{की} \mr{बाळाला} \mr{पडसं} \mr{होतं}.
\end{quote}

\begin{quote}
\emph{havā badallī kī bāḷālā paḍsaṁ hotaṁ.}
\end{quote}

\begin{quote}
When the weather changes the baby gets a runny cold.
\end{quote}

Neuter.

\mr{पडसं} \mr{होणे}, with the dative. Sentences write the neuter as \mr{पडसं}. Close to \mr{सर्दी}.

\begin{cousinweb}
\textbf{\mr{पडसे}} (\emph{paḍse}) is a tadbhava, inherited from Sanskrit \emph{pratiśyāya}, \textquotedblleft{}a cold in the head\textquotedblright{}; CDIAL lists Marathi \emph{paḍsẽ} (CDIAL 8604).
\end{cousinweb}

\begin{grammarlens}[title={What you've built}]
One more word for describing what you lived and felt.
\end{grammarlens}

\subsection*{Your turn}
\begin{itemize}
\raggedright
  \item \emph{Say it:} \emph{paḍse}
  \item \emph{Say it:} \emph{paḍse}, once more
  \item \emph{Say it:} \emph{havā badallī kī bāḷālā paḍsaṁ hotaṁ.} — When the weather changes the baby gets a runny cold.
  \item \emph{Recall:} say the Marathi for weak, then the Marathi for to feel sick, then say \emph{paḍse} again
\end{itemize}

\begin{tcolorbox}[breakable,colback=teal!4,colframe=teal!35!black,title={Before you move on}]
What does \mr{पडसे} mean? (\textquotedblleft{}A runny cold\textquotedblright{}.) Say it once more.
\end{tcolorbox}
\section[\texorpdfstring{\mr{कफ} (kaph)}{कफ (kaph)}]{\mr{कफ} (kaph) — phlegm}

\label{lesson:MR-C481-kaph}



\begin{quote}
Before the new one: say the Marathi for to feel sick, then the Marathi for a runny cold.
\end{quote}

\subsection*{You'll want to know: \mr{कफ}}
\textbf{\mr{कफ}} — \emph{kaph} — \textquotedblleft{}phlegm\textquotedblright{}.

\begin{quote}
\mr{सर्दीमुळे} \mr{छातीत} \mr{कफ} \mr{झाला}.
\end{quote}

\begin{quote}
\emph{sardīmuḷe chātīt kaph jhālā.}
\end{quote}

\begin{quote}
The cold brought phlegm to the chest.
\end{quote}

Masculine.

Masculine; in Ayurveda one of the three humours.

\begin{cousinweb}
\textbf{\mr{कफ}} (\emph{kaph}) is a tatsama, taken from Sanskrit \emph{kapha}, \textquotedblleft{}phlegm\textquotedblright{} (CDIAL 2756).
\end{cousinweb}

\begin{grammarlens}[title={What you've built}]
One more word for describing what you lived and felt.
\end{grammarlens}

\subsection*{Your turn}
\begin{itemize}
\raggedright
  \item \emph{Say it:} \emph{kaph}
  \item \emph{Say it:} \emph{kaph}, once more
  \item \emph{Say it:} \emph{sardīmuḷe chātīt kaph jhālā.} — The cold brought phlegm to the chest.
  \item \emph{Recall:} say the Marathi for to feel sick, then the Marathi for a runny cold, then say \emph{kaph} again
\end{itemize}

\begin{tcolorbox}[breakable,colback=teal!4,colframe=teal!35!black,title={Before you move on}]
What does \mr{कफ} mean? (\textquotedblleft{}Phlegm\textquotedblright{}.) Say it once more.
\end{tcolorbox}
\section[\texorpdfstring{\mr{शिंकणे} (śiṅkṇe)}{शिंकणे (śiṅkṇe)}]{\mr{शिंकणे} (śiṅkṇe) — to sneeze}

\label{lesson:MR-C481-sinkne}



\begin{quote}
Before the new one: say the Marathi for a runny cold, then the Marathi for phlegm.
\end{quote}

\subsection*{You'll want to know: \mr{शिंकणे}}
\textbf{\mr{शिंकणे}} — \emph{śiṅkṇe} — \textquotedblleft{}to sneeze\textquotedblright{}.

\begin{quote}
\mr{धुळीमुळे} \mr{आजोबा} \mr{सारखे} \mr{शिंकत} \mr{होते}.
\end{quote}

\begin{quote}
\emph{dhuḷīmuḷe ājobā sārkhe śiṅkat hote.}
\end{quote}

\begin{quote}
Grandfather kept sneezing because of the dust.
\end{quote}

Past (he …): \mr{तो} \mr{शिंकला}, \emph{to śiṅklā}.

\mr{शिंक}, a sneeze; \mr{शिंक} \mr{येणे}, with the dative.

\begin{cousinweb}
\textbf{\mr{शिंकणे}} (\emph{śiṅkṇe}) is a tadbhava: CDIAL lists Marathi \emph{śikṇẽ}, \textquotedblleft{}sneeze\textquotedblright{}, under Sanskrit \emph{chikkā} (CDIAL 5032).
\end{cousinweb}

\begin{grammarlens}[title={What you've built}]
One more word for describing what you lived and felt.
\end{grammarlens}

\subsection*{Your turn}
\begin{itemize}
\raggedright
  \item \emph{Say it:} \emph{śiṅkṇe}
  \item \emph{Say it:} \emph{śiṅkṇe}, once more
  \item \emph{Say it:} \emph{dhuḷīmuḷe ājobā sārkhe śiṅkat hote.} — Grandfather kept sneezing because of the dust.
  \item \emph{Recall:} say the Marathi for a runny cold, then the Marathi for phlegm, then say \emph{śiṅkṇe} again
\end{itemize}

\begin{tcolorbox}[breakable,colback=teal!4,colframe=teal!35!black,title={Before you move on}]
What does \mr{शिंकणे} mean? (\textquotedblleft{}To sneeze\textquotedblright{}.) Say it once more.
\end{tcolorbox}
\section[\texorpdfstring{\mr{शिंकरणे} (śiṅkarṇe)}{शिंकरणे (śiṅkarṇe)}]{\mr{शिंकरणे} (śiṅkarṇe) — to blow one's nose}

\label{lesson:MR-C481-sinkarne}



\begin{quote}
Before the new one: say the Marathi for phlegm, then the Marathi for to sneeze.
\end{quote}

\subsection*{You'll want to know: \mr{शिंकरणे}}
\textbf{\mr{शिंकरणे}} — \emph{śiṅkarṇe} — \textquotedblleft{}to blow one's nose\textquotedblright{}.

\begin{quote}
\mr{त्याने} \mr{रुमालात} \mr{नाक} \mr{शिंकरलं}.
\end{quote}

\begin{quote}
\emph{tyāne rumālāt nāk śiṅkarlaṁ.}
\end{quote}

\begin{quote}
He blew his nose into a handkerchief.
\end{quote}

Past, with -\mr{ने} (he …): \mr{त्याने} \mr{नाक} \mr{शिंकरलं}, \emph{tyāne nāk śiṅkarlaṁ}.

The object is \mr{नाक}: \mr{नाक} \mr{शिंकरणे}.

\begin{cousinweb}
\textbf{\mr{शिंकरणे}} (\emph{śiṅkarṇe}) is a Marathi verb whose origin is not settled.
\end{cousinweb}

\begin{grammarlens}[title={What you've built}]
One more word for describing what you lived and felt.
\end{grammarlens}

\subsection*{Your turn}
\begin{itemize}
\raggedright
  \item \emph{Say it:} \emph{śiṅkarṇe}
  \item \emph{Say it:} \emph{śiṅkarṇe}, once more
  \item \emph{Say it:} \emph{tyāne rumālāt nāk śiṅkarlaṁ.} — He blew his nose into a handkerchief.
  \item \emph{Recall:} say the Marathi for phlegm, then the Marathi for to sneeze, then say \emph{śiṅkarṇe} again
\end{itemize}

\begin{tcolorbox}[breakable,colback=teal!4,colframe=teal!35!black,title={Before you move on}]
What does \mr{शिंकरणे} mean? (\textquotedblleft{}To blow one's nose\textquotedblright{}.) Say it once more.
\end{tcolorbox}
\section[\texorpdfstring{\mr{खाकरणे} (khākarṇe)}{खाकरणे (khākarṇe)}]{\mr{खाकरणे} (khākarṇe) — to clear one's throat}

\label{lesson:MR-C481-khakarne}



\begin{quote}
Before the new one: say the Marathi for to sneeze, then the Marathi for to blow one's nose.
\end{quote}

\subsection*{You'll want to know: \mr{खाकरणे}}
\textbf{\mr{खाकरणे}} — \emph{khākarṇe} — \textquotedblleft{}to clear one's throat\textquotedblright{}.

\begin{quote}
\mr{बोलण्यापूर्वी} \mr{आजोबा} \mr{खाकरले}.
\end{quote}

\begin{quote}
\emph{bolṇyāpūrvī ājobā khākarle.}
\end{quote}

\begin{quote}
Grandfather cleared his throat before speaking.
\end{quote}

Past (he …): \mr{तो} \mr{खाकरला}, \emph{to khākarlā}.

A small cough to clear the throat, or to catch someone's attention.

\begin{cousinweb}
\textbf{\mr{खाकरणे}} (\emph{khākarṇe}) is inherited: CDIAL lists Marathi \emph{khākarṇẽ}, \textquotedblleft{}clearing the throat\textquotedblright{}, under the reconstructed \emph{*khāṭkāra} (CDIAL 3859).
\end{cousinweb}

\begin{grammarlens}[title={What you've built}]
One more word for describing what you lived and felt.
\end{grammarlens}

\subsection*{Your turn}
\begin{itemize}
\raggedright
  \item \emph{Say it:} \emph{khākarṇe}
  \item \emph{Say it:} \emph{khākarṇe}, once more
  \item \emph{Say it:} \emph{bolṇyāpūrvī ājobā khākarle.} — Grandfather cleared his throat before speaking.
  \item \emph{Recall:} say the Marathi for to sneeze, then the Marathi for to blow one's nose, then say \emph{khākarṇe} again
\end{itemize}

\begin{tcolorbox}[breakable,colback=teal!4,colframe=teal!35!black,title={Before you move on}]
What does \mr{खाकरणे} mean? (\textquotedblleft{}To clear one's throat\textquotedblright{}.) Say it once more.
\end{tcolorbox}
\section[(dialogue)]{Looking back: from Tidy and Scatterbrained to Coughs and Colds}

\label{lesson:MR-R481-review-474-481}



\begin{quote}
Say the Marathi for tidy and sloppy.
\end{quote}

\subsection*{Your turn}
Say these aloud:

\begin{itemize}
\raggedright
  \item tidy, sloppy, scatterbrained, a mess, discipline
  \item adventurous, chicken-hearted, carefree, unconcerned, to be scared of
  \item the eyeball, a nostril, a gum, a beard, a moustache
  \item the palm, a fist, the thigh, the calf, the ankle
  \item the skull, the spine, a rib, the gut, a joint
  \item the pulse, a vein, the gullet, a cell, the innermost heart
  \item one's health, a symptom, fatigue, weak, to feel sick
  \item a runny cold, phlegm, to sneeze, to blow one's nose, to clear one's throat
\end{itemize}

\begin{tcolorbox}[breakable,colback=teal!4,colframe=teal!35!black,title={Before you move on}]
Tidy? (\textbf{\mr{नीटनेटका}}, \emph{nīṭneṭkā}.) Sloppy? (\textbf{\mr{गबाळा}}, \emph{gabāḷā}.) Scatterbrained? (\textbf{\mr{वेंधळा}}, \emph{vendhḷā}.) A mess? (\textbf{\mr{पसारा}}, \emph{pasārā}.) Discipline? (\textbf{\mr{शिस्त}}, \emph{śista}.) Adventurous? (\textbf{\mr{साहसी}}, \emph{sāhsī}.) Chicken-hearted? (\textbf{\mr{डरपोक}}, \emph{ḍarpok}.) Carefree? (\textbf{\mr{बिनधास्त}}, \emph{bindhāsta}.) Unconcerned? (\textbf{\mr{बेफिकीर}}, \emph{bephikīr}.) To be scared of? (\textbf{\mr{भिणे}}, \emph{bhiṇe}.) The eyeball? (\textbf{\mr{बुबुळ}}, \emph{bubuḷ}.) A nostril? (\textbf{\mr{नाकपुडी}}, \emph{nākpuḍī}.) A gum? (\textbf{\mr{हिरडी}}, \emph{hirḍī}.) A beard? (\textbf{\mr{दाढी}}, \emph{dāḍhī}.) A moustache? (\textbf{\mr{मिशी}}, \emph{miśī}.) The palm? (\textbf{\mr{तळहात}}, \emph{taḷhāt}.) A fist? (\textbf{\mr{मूठ}}, \emph{mūṭh}.) The thigh? (\textbf{\mr{मांडी}}, \emph{māṇḍī}.) The calf? (\textbf{\mr{पोटरी}}, \emph{poṭrī}.) The ankle? (\textbf{\mr{घोटा}}, \emph{ghoṭā}.) The skull? (\textbf{\mr{कवटी}}, \emph{kavṭī}.) The spine? (\textbf{\mr{कणा}}, \emph{kaṇā}.) A rib? (\textbf{\mr{बरगडी}}, \emph{bargaḍī}.) The gut? (\textbf{\mr{आतडे}}, \emph{ātḍe}.) A joint? (\textbf{\mr{सांधा}}, \emph{sāndhā}.) The pulse? (\textbf{\mr{नाडी}}, \emph{nāḍī}.) A vein? (\textbf{\mr{नस}}, \emph{nas}.) The gullet? (\textbf{\mr{घसा}}, \emph{ghasā}.) A cell? (\textbf{\mr{पेशी}}, \emph{peśī}.) The innermost heart? (\textbf{\mr{काळीज}}, \emph{kāḷīj}.) One's health? (\textbf{\mr{तब्येत}}, \emph{tabyet}.) A symptom? (\textbf{\mr{लक्षण}}, \emph{lakṣaṇ}.) Fatigue? (\textbf{\mr{थकवा}}, \emph{thakvā}.) Weak? (\textbf{\mr{अशक्त}}, \emph{aśakta}.) To feel sick? (\textbf{\mr{मळमळणे}}, \emph{maḷmaḷṇe}.) A runny cold? (\textbf{\mr{पडसे}}, \emph{paḍse}.) Phlegm? (\textbf{\mr{कफ}}, \emph{kaph}.) To sneeze? (\textbf{\mr{शिंकणे}}, \emph{śiṅkṇe}.) To blow one's nose? (\textbf{\mr{शिंकरणे}}, \emph{śiṅkarṇe}.) To clear one's throat? (\textbf{\mr{खाकरणे}}, \emph{khākarṇe}.)
\end{tcolorbox}
